\frontchapter{Abstract}

Lung nodules are followed up by measuring how fast they grow, but computed tomography (CT), the standard tool, gives the patient a radiation dose and a cost at every follow-up, while low-dose digital chest tomosynthesis (DTS) cannot resolve depth because its \ang{30} sweep leaves a ``missing wedge'' of unmeasured information. MedLift-3D is a click-to-3D system that turns one click on a DTS slice into a 3D nodule measurement with a calibrated uncertainty. We first show from projection moments that every DTS view measures a nodule's mass exactly, while its depth extent, which volume needs, is about 15 times worse conditioned; mass and mass doubling time are therefore the measurements DTS can support. The pipeline simulates low-dose 15-view DTS for 583 LIDC-IDRI patients (1,144 radiologist-annotated nodules), reconstructs a 64 mm prism through the chest around the click, fills the missing wedge with a 17.4-million-parameter conditional 3D diffusion model kept consistent with the X-rays, represents the nodule with 3D Gaussians that give its mass in closed form, and wraps every number in conformal prediction intervals and calibrated size-category probabilities. A learned 3D detector proposes where to click and a seen/unseen (null-space) analysis elucidates on which part of every measurement originates from the X-rays and what is learned from the prior. The system components performed as expected: the PSNR of the nodule-region images on the held out data went from 11.88 dB to 15.23 dB after diffusion prior; the CPM of the detector was 0.298, and the baseline (training-free) was 0.017; the conformal intervals achieved 74\% coverage for the test nodules; and the calibrated category probabilities reduced the Brier score for a volume of at least 100 mm$^3$ from 0.54 to 0.17. The null space analysis validated that the X-ray measured nearly all of the mass of a nodule (0.3\% unseen) and that the prior's unseen content is correlated with the truth (skill 0.22). However, the measurements did not achieve the desired target; it was not more precise than volumetry, the volume error of the entire method (MdAPE 94.6\%) was not smaller than the physics-only reconstruction (87.7 -- 91.2\%) and the valid ranges were not narrow enough for use. These failures we attribute to a forward that reduces small nodules, to the mass measurements on the image plane and not from the projections, and to the fact that measurements made in the calibration set were flawed, and demonstrate that better image quality did not necessarily lead to better measurement.

\vspace{1em}\noindent\textbf{Keywords:} digital tomosynthesis, lung nodule volumetry, diffusion models, 3D Gaussians, conformal prediction, explainable AI
